% =================================================================
% FILE: sections/3.Detailed_Design.tex
% Chapter 3 of the sample Report 4: Detailed Design.
% Each subsection is one feature and holds
%   x.y.1 Class Diagram   (one \kfig per module)
%   x.y.2 ... UC-xx.yy    (one \kfig per sequence diagram)
% =================================================================
\section{Detailed Design}

This chapter presents the detailed design of every feature of the system. Each section
(3.x) covers one feature and contains the class diagram of the feature (3.x.1) followed by
its sequence diagrams (3.x.2, 3.x.3, \ldots), one for each use case of the feature. The
class specification of every class that appears in these diagrams --- attributes, methods
and their types --- is tabulated in Chapter~\ref{sec:class-specification} rather than repeated under each
feature. Where several functions are built from the same classes and follow the same
interaction pattern, the diagrams are presented once for the representative function and
the other functions refer to that section.

% =================================================================
% FILE: sections/kvy/sprint1.tex -- Sprint 1 (Community Forum)
% Features are listed in the approved feature order. PNG diagrams are
% rendered next to the .svg exports in images/sprint1/kvy/.
% =================================================================

% ================= View Forum Category List =================
\subsection{View Forum Category List}

\textbf{a.~Class Diagram}

\kfig{0.92\textwidth}{images/sprint1/kvy/classDiagram/View Forum Category List}{Class diagram of View Forum Category List}{fig:cls-view-forum-category-list}

\textbf{b.~Sequence Diagram(s)}

\kfig{0.78\textwidth}{images/sprint1/kvy/sequences/ViewForumCategories}{Sequence diagram of View Forum Category List}{fig:seq-view-forum-category-list}

% ================= View Forum Category Detail =================
\subsection{View Forum Category Detail}

\textbf{a.~Class Diagram}

\kfig{0.92\textwidth}{images/sprint1/kvy/classDiagram/View Forum Category Detail}{Class diagram of View Forum Category Detail}{fig:cls-view-forum-category-detail}

\textbf{b.~Sequence Diagram(s)}

\kfig{0.78\textwidth}{images/sprint1/kvy/sequences/ViewForumCategoryDetail}{Sequence diagram of View Forum Category Detail}{fig:seq-view-forum-category-detail}

% ================= View Post List =================
\subsection{View Post List}

\textbf{a.~Class Diagram}

\kfig{0.92\textwidth}{images/sprint1/kvy/classDiagram/View Post List}{Class diagram of View Post List}{fig:cls-view-post-list}

\textbf{b.~Sequence Diagram(s)}

\kfig{0.78\textwidth}{images/sprint1/kvy/sequences/PostList}{Sequence diagram of View Post List}{fig:seq-view-post-list}

% ================= View Post Detail =================
\subsection{View Post Detail}

\textbf{a.~Class Diagram}

\kfig{0.92\textwidth}{images/sprint1/kvy/classDiagram/View Post Detail}{Class diagram of View Post Detail}{fig:cls-view-post-detail}

\textbf{b.~Sequence Diagram(s)}

\kfig{0.78\textwidth}{images/sprint1/kvy/sequences/PostDetail}{Sequence diagram of View Post Detail}{fig:seq-view-post-detail}

% ================= Create Post =================
\subsection{Create Post}

\textbf{a.~Class Diagram}

\kfig{0.92\textwidth}{images/sprint1/kvy/classDiagram/Create Post}{Class diagram of Create Post}{fig:cls-create-post}

\textbf{b.~Sequence Diagram(s)}

\kfig{0.78\textwidth}{images/sprint1/kvy/sequences/CreatePost}{Sequence diagram of Create Post}{fig:seq-create-post}

% ================= Update Post =================
\subsection{Update Post}

\textbf{a.~Class Diagram}

\kfig{0.92\textwidth}{images/sprint1/kvy/classDiagram/Update Post}{Class diagram of Update Post}{fig:cls-update-post}

\textbf{b.~Sequence Diagram(s)}

\kfig{0.78\textwidth}{images/sprint1/kvy/sequences/EditPost}{Sequence diagram of Update Post}{fig:seq-update-post}

% ================= Delete Post =================
\subsection{Delete Post}

\textbf{a.~Class Diagram}

\kfig{0.92\textwidth}{images/sprint1/kvy/classDiagram/Delete Post}{Class diagram of Delete Post}{fig:cls-delete-post}

\textbf{b.~Sequence Diagram(s)}

\kfig{0.78\textwidth}{images/sprint1/kvy/sequences/DeletePost}{Sequence diagram of Delete Post}{fig:seq-delete-post}

% ================= Like Post =================
\subsection{Like Post}

\textbf{a.~Class Diagram}

\kfig{0.92\textwidth}{images/sprint1/kvy/classDiagram/Like Post}{Class diagram of Like Post}{fig:cls-like-post}

\textbf{b.~Sequence Diagram(s)}

\kfig{0.78\textwidth}{images/sprint1/kvy/sequences/LikePost}{Sequence diagram of Like Post}{fig:seq-like-post}

% ================= Unlike Post =================
\subsection{Unlike Post}

\textbf{a.~Class Diagram}

\kfig{0.92\textwidth}{images/sprint1/kvy/classDiagram/Unlike Post}{Class diagram of Unlike Post}{fig:cls-unlike-post}

\textbf{b.~Sequence Diagram(s)}

\kfig{0.78\textwidth}{images/sprint1/kvy/sequences/UnlikePost}{Sequence diagram of Unlike Post}{fig:seq-unlike-post}

% ================= View Post Likes =================
\subsection{View Post Likes}

\textbf{a.~Class Diagram}

\kfig{0.92\textwidth}{images/sprint1/kvy/classDiagram/View Post Likes}{Class diagram of View Post Likes}{fig:cls-view-post-likes}

\textbf{b.~Sequence Diagram(s)}

\kfig{0.78\textwidth}{images/sprint1/kvy/sequences/ViewLikePostDetail}{Sequence diagram of View Post Likes}{fig:seq-view-post-likes}

% ================= View Post Comments =================
\subsection{View Post Comments}

\textbf{a.~Class Diagram}

\kfig{0.92\textwidth}{images/sprint1/kvy/classDiagram/View Post Comment}{Class diagram of View Post Comments}{fig:cls-view-post-comments}

\textbf{b.~Sequence Diagram(s)}

\kfig{0.78\textwidth}{images/sprint1/kvy/sequences/PostComments}{Sequence diagram of View Post Comments}{fig:seq-view-post-comments}

% ================= Create Comment =================
\subsection{Create Comment}

\textbf{a.~Class Diagram}

\kfig{0.92\textwidth}{images/sprint1/kvy/classDiagram/Create Comment}{Class diagram of Create Comment}{fig:cls-create-comment}

\textbf{b.~Sequence Diagram(s)}

\kfig{0.78\textwidth}{images/sprint1/kvy/sequences/AddComment}{Sequence diagram of Create Comment}{fig:seq-create-comment}

% ================= Update Comment =================
\subsection{Update Comment}

\textbf{a.~Class Diagram}

\kfig{0.92\textwidth}{images/sprint1/kvy/classDiagram/Update Comment}{Class diagram of Update Comment}{fig:cls-update-comment}

\textbf{b.~Sequence Diagram(s)}

\kfig{0.78\textwidth}{images/sprint1/kvy/sequences/EditComment}{Sequence diagram of Update Comment}{fig:seq-update-comment}

% ================= Delete Comment =================
\subsection{Delete Comment}

\textbf{a.~Class Diagram}

\kfig{0.92\textwidth}{images/sprint1/kvy/classDiagram/Delete Comment}{Class diagram of Delete Comment}{fig:cls-delete-comment}

\textbf{b.~Sequence Diagram(s)}

\kfig{0.78\textwidth}{images/sprint1/kvy/sequences/DeleteComment}{Sequence diagram of Delete Comment}{fig:seq-delete-comment}

% ================= Like Comment =================
\subsection{Like Comment}

\textbf{a.~Class Diagram}

\kfig{0.92\textwidth}{images/sprint1/kvy/classDiagram/Like Comment}{Class diagram of Like Comment}{fig:cls-like-comment}

\textbf{b.~Sequence Diagram(s)}

\kfig{0.78\textwidth}{images/sprint1/kvy/sequences/LikeComment}{Sequence diagram of Like Comment}{fig:seq-like-comment}

% ================= Unlike Comment =================
\subsection{Unlike Comment}

\textbf{a.~Class Diagram}

\kfig{0.92\textwidth}{images/sprint1/kvy/classDiagram/Unlike Comment}{Class diagram of Unlike Comment}{fig:cls-unlike-comment}

\textbf{b.~Sequence Diagram(s)}

\kfig{0.78\textwidth}{images/sprint1/kvy/sequences/UnlikeComment}{Sequence diagram of Unlike Comment}{fig:seq-unlike-comment}

% ================= Report Content =================
\subsection{Report Content}

\textbf{a.~Class Diagram}

\kfig{0.92\textwidth}{images/sprint1/kvy/classDiagram/Report Content}{Class diagram of Report Content}{fig:cls-report-content}

\textbf{b.~Sequence Diagram(s)}

\kfig{0.78\textwidth}{images/sprint1/kvy/sequences/ReportContent}{Sequence diagram of Report Content}{fig:seq-report-content}


% \input{sections/nhuy/sprint1}
% \input{sections/ntin/sprint1}
% =================================================================
% FILE: sections/ptai/2/sprint1.tex -- Sprint 1
% Job Browsing, Job Management, Bookmarked Jobs Management
% =================================================================

% ================= Job Browsing =================
\subsection{Job Browsing}

Job Browsing allows guests and registered users to search, filter, and view the list of available jobs on the platform, as well as view comprehensive job details.

\subsubsection{Class Diagram}

\kfig{0.92\textwidth}{images/sprint1/ptai/class-diagrams/UC-10_Job_Browsing_Class_Diagram}{Class diagram of Job Browsing}{fig:ptai-cls-job-browsing}

\subsubsection{UC-10.01 -- View Job List}
\kfig{0.78\textwidth}{images/sprint1/ptai/UC-10.01_-_View_Job_List}{Sequence diagram of UC-10.01 View Job List}{fig:ptai-seq-uc-10-01}

\subsubsection{UC-10.02 -- View Job Detail}
\kfig{0.78\textwidth}{images/sprint1/ptai/UC-10.02_-_View_Job_Detail}{Sequence diagram of UC-10.02 View Job Detail}{fig:ptai-seq-uc-10-02}

% ================= Job Management =================
\subsection{Job Management}

Job Management enables clients to manage their posted jobs throughout their lifecycle, including creating new jobs, editing job parameters, publishing drafts, closing completed or cancelled jobs, and soft-deleting job postings.

\subsubsection{Class Diagram}

\kfig{0.92\textwidth}{images/sprint1/ptai/class-diagrams/UC-11_Job_Management_Class_Diagram}{Class diagram of Job Management}{fig:ptai-cls-job-management}

\subsubsection{UC-11.01 -- View My Jobs}
\kfig{0.78\textwidth}{images/sprint1/ptai/UC-11.01_-_View_My_Jobs}{Sequence diagram of UC-11.01 View My Jobs}{fig:ptai-seq-uc-11-01}

\subsubsection{UC-11.02 -- Create Job}
\kfig{0.78\textwidth}{images/sprint1/ptai/UC-11.02_-_Create_Job}{Sequence diagram of UC-11.02 Create Job}{fig:ptai-seq-uc-11-02}

\subsubsection{UC-11.03 -- Publish Job}
\kfig{0.78\textwidth}{images/sprint1/ptai/UC-11.03_-_Publish_Job}{Sequence diagram of UC-11.03 Publish Job}{fig:ptai-seq-uc-11-03}

\subsubsection{UC-11.04 -- Update Job}
\kfig{0.78\textwidth}{images/sprint1/ptai/UC-11.04_-_Update_Job}{Sequence diagram of UC-11.04 Update Job}{fig:ptai-seq-uc-11-04}

\subsubsection{UC-11.05 -- Delete Job}
\kfig{0.78\textwidth}{images/sprint1/ptai/UC-11.05_-_Delete_Job}{Sequence diagram of UC-11.05 Delete Job}{fig:ptai-seq-uc-11-05}

\subsubsection{UC-11.06 -- Change Job Status}
\kfig{0.78\textwidth}{images/sprint1/ptai/UC-11.06_-_Change_Job_Status}{Sequence diagram of UC-11.06 Change Job Status}{fig:ptai-seq-uc-11-06}

% ================= Bookmarked Jobs Management =================
\subsection{Bookmarked Jobs Management}

Bookmarked Jobs Management allows freelancers to save interesting job postings for later review, inspect their list of bookmarked jobs, and remove bookmarks when no longer needed.

\subsubsection{Class Diagram}

\kfig{0.92\textwidth}{images/sprint1/ptai/class-diagrams/UC-12_Bookmarked_Jobs_Management_Class_Diagram}{Class diagram of Bookmarked Jobs Management}{fig:ptai-cls-bookmarked-jobs}

\subsubsection{UC-12.01 -- View Bookmarked Job List}
\kfig{0.78\textwidth}{images/sprint1/ptai/UC-12.01_-_View_Bookmarked_Job_List}{Sequence diagram of UC-12.01 View Bookmarked Job List}{fig:ptai-seq-uc-12-01}

\subsubsection{UC-12.02 -- Bookmark Job}
\kfig{0.78\textwidth}{images/sprint1/ptai/UC-12.02_-_Bookmark_Job}{Sequence diagram of UC-12.02 Bookmark Job}{fig:ptai-seq-uc-12-02}

\subsubsection{UC-12.03 -- Remove Bookmark}
\kfig{0.78\textwidth}{images/sprint1/ptai/UC-12.03_-_Remove_Bookmark}{Sequence diagram of UC-12.03 Remove Bookmark}{fig:ptai-seq-uc-12-03}

% =================================================================
% FILE: sections/tthuan/2/sprint1.tex -- Sprint 1
% Authentication and Role Management
% =================================================================

% ================= Authentication =================
\subsection{Authentication}

The authentication module covers account registration, email/password and
Google OAuth login, logout, password recovery, and the JWT/session flow shared
by all protected endpoints.

\subsubsection{Class Diagram}

\kfig{0.92\textwidth}{images/sprint1/tthuan/classDiagram/AuthClasses}{Class diagram of Authentication}{fig:tthuan-cls-auth}

\subsubsection{UC-01 -- Login with Google}
\kfig{0.78\textwidth}{images/sprint1/tthuan/sequences/LoginWithGoogle}{Sequence diagram of UC-01 Login with Google}{fig:tthuan-seq-login-google}

\subsubsection{UC-02 -- Login}
\kfig{0.78\textwidth}{images/sprint1/tthuan/sequences/Login}{Sequence diagram of UC-02 Login}{fig:tthuan-seq-login}

\subsubsection{UC-03 -- Logout}
\kfig{0.78\textwidth}{images/sprint1/tthuan/sequences/Logout}{Sequence diagram of UC-03 Logout}{fig:tthuan-seq-logout}

\subsubsection{UC-04 -- Register}
\kfig{0.78\textwidth}{images/sprint1/tthuan/sequences/Register}{Sequence diagram of UC-04 Register}{fig:tthuan-seq-register}

\subsubsection{UC-05 -- Forgot Password}
\kfig{0.78\textwidth}{images/sprint1/tthuan/sequences/ForgotPassword}{Sequence diagram of UC-05 Forgot Password}{fig:tthuan-seq-forgot-password}

% ================= Role Management =================
\subsection{Role Management}

Role Management lets an administrator inspect the role catalogue, maintain
role metadata, and replace the permission set assigned to a role.

\subsubsection{Class Diagram}

\kfig{0.92\textwidth}{images/sprint1/tthuan/classDiagram/RoleManagement}{Class diagram of Role Management}{fig:tthuan-cls-role-management}

\subsubsection{UC-06.01 -- View Role List}
\kfig{0.78\textwidth}{images/sprint1/tthuan/sequences/ViewRoleList}{Sequence diagram of UC-06.01 View Role List}{fig:tthuan-seq-view-role-list}

\subsubsection{UC-06.02 -- View Role Detail}
\kfig{0.78\textwidth}{images/sprint1/tthuan/sequences/ViewRoleDetail}{Sequence diagram of UC-06.02 View Role Detail}{fig:tthuan-seq-view-role-detail}

\subsubsection{UC-06.03 -- Create Role}
\kfig{0.78\textwidth}{images/sprint1/tthuan/sequences/CreateRole}{Sequence diagram of UC-06.03 Create Role}{fig:tthuan-seq-create-role}

\subsubsection{UC-06.04 -- Update Role}
\kfig{0.78\textwidth}{images/sprint1/tthuan/sequences/UpdateRole}{Sequence diagram of UC-06.04 Update Role}{fig:tthuan-seq-update-role}

\subsubsection{UC-06.05 -- Delete Role}
\kfig{0.78\textwidth}{images/sprint1/tthuan/sequences/DeleteRole}{Sequence diagram of UC-06.05 Delete Role}{fig:tthuan-seq-delete-role}

\subsubsection{UC-06.06 -- Update Role Permissions}
\kfig{0.78\textwidth}{images/sprint1/tthuan/sequences/UpdateRolePermissions}{Sequence diagram of UC-06.06 Update Role Permissions}{fig:tthuan-seq-update-role-permissions}


% =================================================================
% FILE: sections/kvy/2/sprint2.tex -- Sprint 2 (Conversation)
% UC-27 Manage Conversation. Figures are the vector PDFs exported from
% class-diagrams/UC-27_Manage_Conversation_Class_Diagram.drawio and
% sequence-diagrams/UC-27.<nn>-*.drawio into images/sprint2/kvy/.
% =================================================================

% ================= Manage Conversation =================
\subsection{Manage Conversation}

Manage Conversation covers the UC-27 private messaging module: listing and
starting one-to-one conversations, opening a conversation, sending text
messages and files, marking messages as read, recalling a message, and
pinning or unpinning a conversation. Delivery is served both by the REST
endpoints and by the realtime gateway in the same namespace.

\subsubsection{Class Diagram}

\kfig{0.92\textwidth}{images/sprint2/kvy/classDiagram/ManageConversation}{Class diagram of Manage Conversation}{fig:kvy-cls-manage-conversation}

\subsubsection{UC-27.01 -- View Conversation List}
\kfig{0.78\textwidth}{images/sprint2/kvy/sequences/UC-27.01-view-conversation-list}{Sequence diagram of UC-27.01 View Conversation List}{fig:kvy-seq-uc-27-01}

\subsubsection{UC-27.02 -- Start Conversation}
\kfig{0.78\textwidth}{images/sprint2/kvy/sequences/UC-27.02-start-conversation}{Sequence diagram of UC-27.02 Start Conversation}{fig:kvy-seq-uc-27-02}

\subsubsection{UC-27.03 -- View Conversation Detail}
\kfig{0.78\textwidth}{images/sprint2/kvy/sequences/UC-27.03-view-conversation-detail}{Sequence diagram of UC-27.03 View Conversation Detail}{fig:kvy-seq-uc-27-03}

\subsubsection{UC-27.04 -- Send Message}
\kfig{0.78\textwidth}{images/sprint2/kvy/sequences/UC-27.04-send-message}{Sequence diagram of UC-27.04 Send Message}{fig:kvy-seq-uc-27-04}

\subsubsection{UC-27.05 -- Mark Message as Read}
\kfig{0.78\textwidth}{images/sprint2/kvy/sequences/UC-27.05-mark-message-as-read}{Sequence diagram of UC-27.05 Mark Message as Read}{fig:kvy-seq-uc-27-05}

\subsubsection{UC-27.06 -- Delete/Recall Message}
\kfig{0.78\textwidth}{images/sprint2/kvy/sequences/UC-27.06-delete-recall-message}{Sequence diagram of UC-27.06 Delete/Recall Message}{fig:kvy-seq-uc-27-06}

\subsubsection{UC-27.07 -- Pin/Unpin Conversation}
\kfig{0.78\textwidth}{images/sprint2/kvy/sequences/UC-27.07-pin-unpin-conversation}{Sequence diagram of UC-27.07 Pin/Unpin Conversation}{fig:kvy-seq-uc-27-07}

\subsubsection{UC-27.08 -- Send File}
\kfig{0.78\textwidth}{images/sprint2/kvy/sequences/UC-27.08-send-file}{Sequence diagram of UC-27.08 Send File}{fig:kvy-seq-uc-27-08}

% \input{sections/nhuy/sprint2}
% \input{sections/ntin/sprint2}
% =================================================================
% FILE: sections/ptai/2/sprint2.tex -- Sprint 2
% Saved Search Management and Job Alert Management
% =================================================================

% ================= Saved Search Management =================
\subsection{Saved Search Management}

Saved Search Management allows freelancers to save complex job filter criteria, view their list of saved searches, inspect search details, update filter parameters, and remove obsolete searches.

\textbf{a.~Class Diagram}

\kfig{0.92\textwidth}{images/sprint2/ptai/class-diagrams/UC-22_Saved_Search_Management_Class_Diagram}{Class diagram of Saved Search Management}{fig:ptai-cls-saved-search}

\textbf{b.~Sequence Diagram(s)}

\subsubsection{UC-22.01 -- View Saved Search List}
\kfig{0.78\textwidth}{images/sprint2/ptai/UC-22.01_-_View_Saved_Search_List}{Sequence diagram of UC-22.01 View Saved Search List}{fig:ptai-seq-uc-22-01}

\subsubsection{UC-22.02 -- View Saved Search Detail}
\kfig{0.78\textwidth}{images/sprint2/ptai/UC-22.02_-_View_Saved_Search_Detail}{Sequence diagram of UC-22.02 View Saved Search Detail}{fig:ptai-seq-uc-22-02}

\subsubsection{UC-22.03 -- Create Saved Search}
\kfig{0.78\textwidth}{images/sprint2/ptai/UC-22.03_-_Create_Saved_Search}{Sequence diagram of UC-22.03 Create Saved Search}{fig:ptai-seq-uc-22-03}

\subsubsection{UC-22.04 -- Update Saved Search}
\kfig{0.78\textwidth}{images/sprint2/ptai/UC-22.04_-_Update_Saved_Search}{Sequence diagram of UC-22.04 Update Saved Search}{fig:ptai-seq-uc-22-04}

\subsubsection{UC-22.05 -- Delete Saved Search}
\kfig{0.78\textwidth}{images/sprint2/ptai/UC-22.05_-_Delete_Saved_Search}{Sequence diagram of UC-22.05 Delete Saved Search}{fig:ptai-seq-uc-22-05}

% ================= Job Alert Management =================
\subsection{Job Alert Management}

Job Alert Management enables freelancers to configure automated notifications based on job keywords, budget ranges, and category criteria, ensuring they receive prompt alerts when matching jobs are published.

\textbf{a.~Class Diagram}

\kfig{0.92\textwidth}{images/sprint2/ptai/class-diagrams/UC-23_Job_Alert_Management_Class_Diagram}{Class diagram of Job Alert Management}{fig:ptai-cls-job-alert}

\textbf{b.~Sequence Diagram(s)}

\subsubsection{UC-23.01 -- View Job Alert List}
\kfig{0.78\textwidth}{images/sprint2/ptai/UC-23.01_-_View_Job_Alert_List}{Sequence diagram of UC-23.01 View Job Alert List}{fig:ptai-seq-uc-23-01}

\subsubsection{UC-23.02 -- View Job Alert Detail}
\kfig{0.78\textwidth}{images/sprint2/ptai/UC-23.02_-_View_Job_Alert_Detail}{Sequence diagram of UC-23.02 View Job Alert Detail}{fig:ptai-seq-uc-23-02}

\subsubsection{UC-23.03 -- Create Job Alert}
\kfig{0.78\textwidth}{images/sprint2/ptai/UC-23.03_-_Create_Job_Alert}{Sequence diagram of UC-23.03 Create Job Alert}{fig:ptai-seq-uc-23-03}

\subsubsection{UC-23.04 -- Update Job Alert}
\kfig{0.78\textwidth}{images/sprint2/ptai/UC-23.04_-_Update_Job_Alert}{Sequence diagram of UC-23.04 Update Job Alert}{fig:ptai-seq-uc-23-04}

\subsubsection{UC-23.05 -- Delete Job Alert}
\kfig{0.78\textwidth}{images/sprint2/ptai/UC-23.05_-_Delete_Job_Alert}{Sequence diagram of UC-23.05 Delete Job Alert}{fig:ptai-seq-uc-23-05}

% =================================================================
% FILE: sections/tthuan/2/sprint2.tex -- Sprint 2
% Permission, Session, and Forum Category Management
% =================================================================

% ================= Permission Management =================
\subsection{Permission Management}

Permission Management exposes the route-level permission catalogue so that an
administrator can inspect the permissions available for role assignment.

\textbf{a.~Class Diagram}

\kfig{0.92\textwidth}{images/sprint2/tthuan/classDiagram/PermissionManagement}{Class diagram of Permission Management}{fig:tthuan-cls-permission-management}

\textbf{b.~Sequence Diagram(s)}

\subsubsection{UC-15.01 -- View Permission List}
\kfig{0.78\textwidth}{images/sprint2/tthuan/sequences/ViewPermissionList}{Sequence diagram of UC-15.01 View Permission List}{fig:tthuan-seq-view-permission-list}

\subsubsection{UC-15.02 -- View Permission Detail}
\kfig{0.78\textwidth}{images/sprint2/tthuan/sequences/ViewPermissionDetail}{Sequence diagram of UC-15.02 View Permission Detail}{fig:tthuan-seq-view-permission-detail}

% ================= Session Management =================
\subsection{Session Management}

Session Management gives an authenticated user visibility over the devices
that currently hold a refresh-token session and allows a session to be revoked
remotely.

\textbf{a.~Class Diagram}

\kfig{0.92\textwidth}{images/sprint2/tthuan/classDiagram/SessionManagement}{Class diagram of Session Management}{fig:tthuan-cls-session-management}

\textbf{b.~Sequence Diagram(s)}

\subsubsection{UC-16.01 -- View Session List}
\kfig{0.78\textwidth}{images/sprint2/tthuan/sequences/ViewSessionList}{Sequence diagram of UC-16.01 View Session List}{fig:tthuan-seq-view-session-list}

\subsubsection{UC-16.02 -- View Session Detail}
\kfig{0.78\textwidth}{images/sprint2/tthuan/sequences/ViewSessionDetail}{Sequence diagram of UC-16.02 View Session Detail}{fig:tthuan-seq-view-session-detail}

\subsubsection{UC-16.03 -- Revoke Session}
\kfig{0.78\textwidth}{images/sprint2/tthuan/sequences/RevokeSession}{Sequence diagram of UC-16.03 Revoke Session}{fig:tthuan-seq-revoke-session}

% ================= Forum Category Management =================
\subsection{Forum Category Management}

Forum Category Management covers the administrator-side category catalogue and
the public category views used by the forum. Deletion is a soft delete, so a
category can be restored later.

\textbf{a.~Class Diagram}

\kfig{0.92\textwidth}{images/sprint2/tthuan/classDiagram/ForumCategoryManagement}{Class diagram of Forum Category Management}{fig:tthuan-cls-forum-category-management}

\textbf{b.~Sequence Diagram(s)}

\subsubsection{UC-17.01 -- Create Forum Category}
\kfig{0.78\textwidth}{images/sprint2/tthuan/sequences/CreateForumCategory}{Sequence diagram of UC-17.01 Create Forum Category}{fig:tthuan-seq-create-forum-category}

\subsubsection{UC-17.02 -- Update Forum Category}
\kfig{0.78\textwidth}{images/sprint2/tthuan/sequences/UpdateForumCategory}{Sequence diagram of UC-17.02 Update Forum Category}{fig:tthuan-seq-update-forum-category}

\subsubsection{UC-17.03 -- Delete Forum Category}
\kfig{0.78\textwidth}{images/sprint2/tthuan/sequences/DeleteForumCategory}{Sequence diagram of UC-17.03 Delete Forum Category}{fig:tthuan-seq-delete-forum-category}

\subsubsection{UC-17.06 -- Restore Forum Category}
\kfig{0.78\textwidth}{images/sprint2/tthuan/sequences/RestoreForumCategory}{Sequence diagram of UC-17.06 Restore Forum Category}{fig:tthuan-seq-restore-forum-category}

\subsubsection{UC-14.01 -- View Forum Category List}
\kfig{0.78\textwidth}{images/sprint2/tthuan/sequences/ViewForumCategoryList}{Sequence diagram of UC-14.01 View Forum Category List}{fig:tthuan-seq-view-forum-category-list}

\subsubsection{UC-14.02 -- View Forum Category Detail}
\kfig{0.78\textwidth}{images/sprint2/tthuan/sequences/ViewForumCategoryDetail}{Sequence diagram of UC-14.02 View Forum Category Detail}{fig:tthuan-seq-view-forum-category-detail}


% =================================================================
% FILE: sections/kvy/2/sprint3.tex -- Sprint 3 (Administration)
% UC-34 Manage Moderation, UC-35 Manage ID Verification and
% UC-36 Manage Forum Trash. Figures are the vector PDFs exported from
% class-diagrams/UC-3[456]_*.drawio and sequence-diagrams/UC-3[456].<nn>-*
% into images/sprint3/kvy/.
% =================================================================

% ================= Manage Moderation =================
\subsection{Manage Moderation}

Manage Moderation covers the UC-34 review workflow of the administration
console: listing the posts waiting for review, opening a flagged post and
approving or rejecting it. A review writes the moderation status, score and
reviewer back to the forum post.

\textbf{a.~Class Diagram}

\kfig{0.92\textwidth}{images/sprint3/kvy/classDiagram/ManageModeration}{Class diagram of Manage Moderation}{fig:kvy-cls-manage-moderation}

\textbf{b.~Sequence Diagram(s)}

\subsubsection{UC-34.01 -- View Moderation Queue}
\kfig{0.78\textwidth}{images/sprint3/kvy/sequences/UC-34.01-view-moderation-queue}{Sequence diagram of UC-34.01 View Moderation Queue}{fig:kvy-seq-uc-34-01}

\subsubsection{UC-34.02 -- View Moderation Item Detail}
\kfig{0.78\textwidth}{images/sprint3/kvy/sequences/UC-34.02-view-moderation-item-detail}{Sequence diagram of UC-34.02 View Moderation Item Detail}{fig:kvy-seq-uc-34-02}

\subsubsection{UC-34.03 -- Resolve Moderation Item}
\kfig{0.78\textwidth}{images/sprint3/kvy/sequences/UC-34.03-resolve-moderation-item}{Sequence diagram of UC-34.03 Resolve Moderation Item}{fig:kvy-seq-uc-34-03}

% ================= Manage ID Verification =================
\subsection{Manage ID Verification}

Manage ID Verification covers the UC-35 identity review workflow: listing the
submitted identity documents, opening a request with its uploaded file, and
approving or rejecting the request. A decision updates the verification status
and the reviewer of the identity document.

\textbf{a.~Class Diagram}

\kfig{0.92\textwidth}{images/sprint3/kvy/classDiagram/ManageIDVerification}{Class diagram of Manage ID Verification}{fig:kvy-cls-manage-id-verification}

\textbf{b.~Sequence Diagram(s)}

\subsubsection{UC-35.01 -- View ID Verification Request List}
\kfig{0.78\textwidth}{images/sprint3/kvy/sequences/UC-35.01-view-id-verification-request-list}{Sequence diagram of UC-35.01 View ID Verification Request List}{fig:kvy-seq-uc-35-01}

\subsubsection{UC-35.02 -- View ID Verification Request Detail}
\kfig{0.78\textwidth}{images/sprint3/kvy/sequences/UC-35.02-view-id-verification-request-detail}{Sequence diagram of UC-35.02 View ID Verification Request Detail}{fig:kvy-seq-uc-35-02}

\subsubsection{UC-35.03 -- Approve ID Verification}
\kfig{0.78\textwidth}{images/sprint3/kvy/sequences/UC-35.03-approve-id-verification}{Sequence diagram of UC-35.03 Approve ID Verification}{fig:kvy-seq-uc-35-03}

\subsubsection{UC-35.04 -- Reject ID Verification}
\kfig{0.78\textwidth}{images/sprint3/kvy/sequences/UC-35.04-reject-id-verification}{Sequence diagram of UC-35.04 Reject ID Verification}{fig:kvy-seq-uc-35-04}

% ================= Manage Forum Trash =================
\subsection{Manage Forum Trash}

Manage Forum Trash covers the UC-36 recovery workflow: listing the soft-deleted
posts and comments kept for review, opening one deleted item, and restoring it
so that it becomes visible in the forum again.

\textbf{a.~Class Diagram}

\kfig{0.92\textwidth}{images/sprint3/kvy/classDiagram/ManageForumTrash}{Class diagram of Manage Forum Trash}{fig:kvy-cls-manage-forum-trash}

\textbf{b.~Sequence Diagram(s)}

\subsubsection{UC-36.01 -- View Deleted Posts/Comments List}
\kfig{0.78\textwidth}{images/sprint3/kvy/sequences/UC-36.01-view-deleted-posts-comments-list}{Sequence diagram of UC-36.01 View Deleted Posts/Comments List}{fig:kvy-seq-uc-36-01}

\subsubsection{UC-36.02 -- View Deleted Posts/Comments Detail}
\kfig{0.78\textwidth}{images/sprint3/kvy/sequences/UC-36.02-view-deleted-posts-comments-detail}{Sequence diagram of UC-36.02 View Deleted Posts/Comments Detail}{fig:kvy-seq-uc-36-02}

\subsubsection{UC-36.03 -- Restore Deleted Posts/Comments}
\kfig{0.78\textwidth}{images/sprint3/kvy/sequences/UC-36.03-restore-deleted-posts-comments}{Sequence diagram of UC-36.03 Restore Deleted Posts/Comments}{fig:kvy-seq-uc-36-03}

% \input{sections/nhuy/sprint3}
% \input{sections/ntin/sprint3}
% =================================================================
% FILE: sections/ptai/2/sprint3.tex -- Sprint 3
% Proposal Management
% =================================================================

% ================= Proposal Management =================
\subsection{Proposal Management}

Proposal Management covers the full proposal workflow between freelancers and clients. Freelancers can submit proposals with cover letters, bid amounts, and delivery estimates, view their submitted proposals, modify pending proposals, or withdraw them. Clients can view all proposals submitted for their jobs, inspect detailed freelancer bids, and accept or reject proposals to initiate contracting.

\subsubsection{Class Diagram}

\kfig{0.92\textwidth}{images/sprint3/ptai/class-diagrams/UC-32_Freelancer_Proposal_Management_Class_Diagram}{Class diagram of Freelancer Proposal Management}{fig:ptai-cls-freelancer-proposal}

\kfig{0.92\textwidth}{images/sprint3/ptai/class-diagrams/UC-32_Client_Proposal_Management_Class_Diagram}{Class diagram of Client Proposal Management}{fig:ptai-cls-client-proposal}

\subsubsection{UC-32.01 -- View My Proposal List}
\kfig{0.78\textwidth}{images/sprint3/ptai/UC-32.01_-_View_My_Proposal_List}{Sequence diagram of UC-32.01 View My Proposal List}{fig:ptai-seq-uc-32-01}

\subsubsection{UC-32.02 -- View My Proposal Detail}
\kfig{0.78\textwidth}{images/sprint3/ptai/UC-32.02_-_View_My_Proposal_Detail}{Sequence diagram of UC-32.02 View My Proposal Detail}{fig:ptai-seq-uc-32-02}

\subsubsection{UC-32.03 -- Submit Proposal}
\kfig{0.78\textwidth}{images/sprint3/ptai/UC-32.03_-_Submit_Proposal}{Sequence diagram of UC-32.03 Submit Proposal}{fig:ptai-seq-uc-32-03}

\subsubsection{UC-32.04 -- Update Proposal}
\kfig{0.78\textwidth}{images/sprint3/ptai/UC-32.04_-_Update_Proposal}{Sequence diagram of UC-32.04 Update Proposal}{fig:ptai-seq-uc-32-04}

\subsubsection{UC-32.05 -- Withdraw Proposal}
\kfig{0.78\textwidth}{images/sprint3/ptai/UC-32.05_-_Withdraw_Proposal}{Sequence diagram of UC-32.05 Withdraw Proposal}{fig:ptai-seq-uc-32-05}

\subsubsection{UC-32.06 -- Accept Proposal}
\kfig{0.78\textwidth}{images/sprint3/ptai/UC-32.06_-_Accept_Proposal}{Sequence diagram of UC-32.06 Accept Proposal}{fig:ptai-seq-uc-32-06}

\subsubsection{UC-32.07 -- Reject Proposal}
\kfig{0.78\textwidth}{images/sprint3/ptai/UC-32.07_-_Reject_Proposal}{Sequence diagram of UC-32.07 Reject Proposal}{fig:ptai-seq-uc-32-07}

\subsubsection{UC-32.08 -- View Job Proposal List}
\kfig{0.78\textwidth}{images/sprint3/ptai/UC-32.08_-_View_Job_Proposal_List}{Sequence diagram of UC-32.08 View Job Proposal List}{fig:ptai-seq-uc-32-08}

\subsubsection{UC-32.09 -- View Proposal Detail}
\kfig{0.78\textwidth}{images/sprint3/ptai/UC-32.09_-_View_Proposal_Detail}{Sequence diagram of UC-32.09 View Proposal Detail}{fig:ptai-seq-uc-32-09}

% =================================================================
% FILE: sections/tthuan/2/sprint3.tex -- Sprint 3
% User Management
% =================================================================

% ================= User Management =================
\subsection{User Management}

User Management provides the administrator with a searchable user catalogue,
a consolidated account detail view, and dedicated editors for the Client and
Freelancer profiles attached to an account.

\textbf{a.~Class Diagram}

\kfig{0.92\textwidth}{images/sprint3/tthuan/classDiagram/UserManagement}{Class diagram of User Management}{fig:tthuan-cls-user-management}

\textbf{b.~Sequence Diagram(s)}

\subsubsection{UC-28.01 -- View User List}
\kfig{0.78\textwidth}{images/sprint3/tthuan/sequences/ViewUserList}{Sequence diagram of UC-28.01 View User List}{fig:tthuan-seq-view-user-list}

\subsubsection{UC-28.02 -- View User Detail}
\kfig{0.78\textwidth}{images/sprint3/tthuan/sequences/ViewUserDetail}{Sequence diagram of UC-28.02 View User Detail}{fig:tthuan-seq-view-user-detail}

\subsubsection{UC-28.03 -- Create User}
\kfig{0.78\textwidth}{images/sprint3/tthuan/sequences/CreateUser}{Sequence diagram of UC-28.03 Create User}{fig:tthuan-seq-create-user}

\subsubsection{UC-28.04 -- Update User}
\kfig{0.78\textwidth}{images/sprint3/tthuan/sequences/UpdateUser}{Sequence diagram of UC-28.04 Update User}{fig:tthuan-seq-update-user}

\subsubsection{UC-28.05 -- Update Client Profile}
\kfig{0.78\textwidth}{images/sprint3/tthuan/sequences/UpdateClientProfile}{Sequence diagram of UC-28.05 Update Client Profile}{fig:tthuan-seq-update-client-profile}

\subsubsection{UC-28.06 -- Update Freelancer Basic Info}
\kfig{0.78\textwidth}{images/sprint3/tthuan/sequences/UpdateFreelancerBasicInfo}{Sequence diagram of UC-28.06 Update Freelancer Basic Info}{fig:tthuan-seq-update-freelancer-basic-info}

\subsubsection{UC-28.07 -- Update Freelancer Skills}
\kfig{0.78\textwidth}{images/sprint3/tthuan/sequences/UpdateFreelancerSkills}{Sequence diagram of UC-28.07 Update Freelancer Skills}{fig:tthuan-seq-update-freelancer-skills}

\subsubsection{UC-28.08 -- Update Freelancer Portfolio}
\kfig{0.78\textwidth}{images/sprint3/tthuan/sequences/UpdateFreelancerPortfolio}{Sequence diagram of UC-28.08 Update Freelancer Portfolio}{fig:tthuan-seq-update-freelancer-portfolio}


% \input{sections/kvy/sprint4}
% \input{sections/nhuy/sprint4}
% \input{sections/ntin/sprint4}
% \input{sections/ptai/sprint4}
% =================================================================
% FILE: sections/tthuan/2/sprint4.tex -- Sprint 4
% Banner Management and Skill Management
% =================================================================

% ================= Banner Management =================
\subsection{Banner Management}

Banner Management maintains the advertisement banners shown in configured
page positions. The module supports image upload, lifecycle dates, activation
state, and hard deletion.

\subsubsection{Class Diagram}

\kfig{0.92\textwidth}{images/sprint4/tthuan/classDiagram/BannerManagement}{Class diagram of Banner Management}{fig:tthuan-cls-banner-management}

\subsubsection{UC-37.01 -- View Banner List}
\kfig{0.78\textwidth}{images/sprint4/tthuan/sequences/ViewBannerList}{Sequence diagram of UC-37.01 View Banner List}{fig:tthuan-seq-view-banner-list}

\subsubsection{UC-37.02 -- View Banner Detail}
\kfig{0.78\textwidth}{images/sprint4/tthuan/sequences/ViewBannerDetail}{Sequence diagram of UC-37.02 View Banner Detail}{fig:tthuan-seq-view-banner-detail}

\subsubsection{UC-37.03 -- Create Banner}
\kfig{0.78\textwidth}{images/sprint4/tthuan/sequences/CreateBanner}{Sequence diagram of UC-37.03 Create Banner}{fig:tthuan-seq-create-banner}

\subsubsection{UC-37.04 -- Update Banner}
\kfig{0.78\textwidth}{images/sprint4/tthuan/sequences/UpdateBanner}{Sequence diagram of UC-37.04 Update Banner}{fig:tthuan-seq-update-banner}

\subsubsection{UC-37.05 -- Delete Banner}
\kfig{0.78\textwidth}{images/sprint4/tthuan/sequences/DeleteBanner}{Sequence diagram of UC-37.05 Delete Banner}{fig:tthuan-seq-delete-banner}

% ================= Skill Management =================
\subsection{Skill Management}

Skill Management owns the master skill catalogue used by freelancer profiles
and job postings. Skills are soft-deleted so that existing references remain
valid, and a deleted skill can be restored.

\subsubsection{Class Diagram}

\kfig{0.92\textwidth}{images/sprint4/tthuan/classDiagram/SkillManagement}{Class diagram of Skill Management}{fig:tthuan-cls-skill-management}

\subsubsection{UC-38.01 -- View Skill List}
\kfig{0.78\textwidth}{images/sprint4/tthuan/sequences/ViewSkillList}{Sequence diagram of UC-38.01 View Skill List}{fig:tthuan-seq-view-skill-list}

\subsubsection{UC-38.02 -- View Skill Detail}
\kfig{0.78\textwidth}{images/sprint4/tthuan/sequences/ViewSkillDetail}{Sequence diagram of UC-38.02 View Skill Detail}{fig:tthuan-seq-view-skill-detail}

\subsubsection{UC-38.03 -- Create Skill}
\kfig{0.78\textwidth}{images/sprint4/tthuan/sequences/CreateSkill}{Sequence diagram of UC-38.03 Create Skill}{fig:tthuan-seq-create-skill}

\subsubsection{UC-38.04 -- Update Skill}
\kfig{0.78\textwidth}{images/sprint4/tthuan/sequences/UpdateSkill}{Sequence diagram of UC-38.04 Update Skill}{fig:tthuan-seq-update-skill}

\subsubsection{UC-38.05 -- Delete Skill}
\kfig{0.78\textwidth}{images/sprint4/tthuan/sequences/DeleteSkill}{Sequence diagram of UC-38.05 Delete Skill}{fig:tthuan-seq-delete-skill}

\subsubsection{UC-38.06 -- Restore Skill}
\kfig{0.78\textwidth}{images/sprint4/tthuan/sequences/RestoreSkill}{Sequence diagram of UC-38.06 Restore Skill}{fig:tthuan-seq-restore-skill}


% \input{sections/nhuy/sprint5}
% \input{sections/ntin/sprint5}
% \input{sections/ptai/sprint5}
% =================================================================
% FILE: sections/tthuan/2/sprint5.tex -- Sprint 5
% Expert Profile and Job Category Management
% =================================================================

% ================= Expert Profile Management =================
\subsection{Expert Profile Management}

The administrator-side Expert Profile editor updates an expert account
directly. Any pending self-service revision is superseded, and the account
owner is notified of the administrative change.

\textbf{a.~Class Diagram}

\kfig{0.92\textwidth}{images/sprint5/tthuan/classDiagram/ExpertProfileManagement}{Class diagram of Expert Profile Management}{fig:tthuan-cls-expert-profile-management}

\textbf{b.~Sequence Diagram(s)}

\subsubsection{UC-28.09 -- Update Expert Profile}
\kfig{0.78\textwidth}{images/sprint5/tthuan/sequences/UpdateExpertProfile}{Sequence diagram of UC-28.09 Update Expert Profile}{fig:tthuan-seq-update-expert-profile}

% ================= Job Category Management =================
\subsection{Job Category Management}

Job Category Management maintains the taxonomy used by job postings. A
category can be created, updated, soft-deleted, and restored; deletion is
blocked while active jobs still reference the category.

\textbf{a.~Class Diagram}

\kfig{0.92\textwidth}{images/sprint5/tthuan/classDiagram/JobCategoryManagement}{Class diagram of Job Category Management}{fig:tthuan-cls-job-category-management}

\textbf{b.~Sequence Diagram(s)}

\subsubsection{UC-46.01 -- View Job Category List}
\kfig{0.78\textwidth}{images/sprint5/tthuan/sequences/ViewJobCategoryList}{Sequence diagram of UC-46.01 View Job Category List}{fig:tthuan-seq-view-job-category-list}

\subsubsection{UC-46.02 -- View Job Category Detail}
\kfig{0.78\textwidth}{images/sprint5/tthuan/sequences/ViewJobCategoryDetail}{Sequence diagram of UC-46.02 View Job Category Detail}{fig:tthuan-seq-view-job-category-detail}

\subsubsection{UC-46.03 -- Create Job Category}
\kfig{0.78\textwidth}{images/sprint5/tthuan/sequences/CreateJobCategory}{Sequence diagram of UC-46.03 Create Job Category}{fig:tthuan-seq-create-job-category}

\subsubsection{UC-46.04 -- Update Job Category}
\kfig{0.78\textwidth}{images/sprint5/tthuan/sequences/UpdateJobCategory}{Sequence diagram of UC-46.04 Update Job Category}{fig:tthuan-seq-update-job-category}

\subsubsection{UC-46.05 -- Delete Job Category}
\kfig{0.78\textwidth}{images/sprint5/tthuan/sequences/DeleteJobCategory}{Sequence diagram of UC-46.05 Delete Job Category}{fig:tthuan-seq-delete-job-category}

